\documentclass[11pt]{article}
\usepackage[margin=1in]{geometry}
\usepackage{amsmath,amssymb,amsthm,hyperref}
\newtheorem{theorem}{Theorem}
\newtheorem{lemma}[theorem]{Lemma}
\newcommand{\MA}{\operatorname{MA}}
\newcommand{\R}{\mathbb R}
\title{A generated radial finite-energy equation on an analytic branch}
\author{Research working paper, continuation 385}
\date{9 October 2026}
\begin{document}
\maketitle
\begin{abstract}
Assuming exactly the bounded analytic calculus stated in continuation 384,
we construct its decreasing-limit radial extension. Its mixed operators
are genuine limits of bounded Monge--Amp\`ere products, have raw mass two,
and satisfy radial comparison and decreasing continuity. Energy sublevels
are compact in a separated topology away from one rigid point. Every
measure in a precisely specified branch family satisfying an explicit
tail-energy condition has a unique radial solution modulo constants.
This is a conditional special case, not the intrinsic all-direction
theory or a solution of the full original P16.7. The bounded foundations
and this new extension remain unreviewed.
\end{abstract}

\section{The exact bounded premises}
Use the actual quotient, ordinary bounded class, reference $u_0$, depth
$T$, branch $W$, Gauss points $\xi_t$ for $t\geq1$, and raw mass two
of continuation 384. All its explicitly numbered bounded calculus
assumptions are retained, including ordinary locality, coefficient
multiplication, uniform product continuity, and the calibrated
degree-one translated-torus rule. Its Proposition 1 supplies, for every
nonnegative compact finite step profile $h\leq1$, nonincreasing and
constant on $[0,1)$,
\[
 u_h=u_0-g_h(T),\qquad g_h(t)=\int_0^t h(v)\,\mathrm dv,
\]
\[
 \MA(u_h)|_W=\xi_*[-\mathrm d(h^2)],\qquad
 E(u_h)=-\frac16-\frac13\int_0^\infty h^3\,\mathrm dt.
\]
These are conditional analytic bounded-product identities, not imported
accepted grants. Continuation 384 proves them by a complete softmax
gradient calculation against analytic tests, not an inverse measure
construction. Its full source was read, including the zero-profile repair.

Put $\mu_0=2\delta_{i(0)}$ and
$\lambda=\MA(u_0+D_1)|_{X\setminus W}$. This is an actual positive
bounded-operator measure of mass one. No particular support point of
$\lambda$ is presumed. For a finite profile with $c=h(0)$ all cutoffs
$D_N$, $N\geq1$, agree with $D_1$ on the ordinary open $\{T<1\}$,
which contains $X\setminus W$. Ordinary locality and the bounded pencil
identity therefore give
\begin{equation}\label{eq:bounded}
 \MA(u_h)=(1-c^2)\mu_0+c^2\lambda+\xi_*[-\mathrm d(h^2)].
\end{equation}
Here and below the Stieltjes measure lives on $[1,\infty)$; profiles are
right continuous on $[1,\infty)$, and a jump at 1 is included.
Indeed the pencil identity follows by bilinearity and
$A\wedge\mathrm d'\mathrm d''D_1=0$ from the exact bounded hypotheses.
Polarizing \eqref{eq:bounded} for two profiles with initial values $c,d$
likewise gives the bounded mixed product
\begin{equation}\label{eq:mixed}
 B(u_h,u_k)=(1-cd)\mu_0+cd\lambda+\xi_*[-\mathrm d(hk)].
\end{equation}
No outside mass is guessed; locality fixes it through $\lambda$.

\section{Construction of the generated extension}
Let $\mathcal H$ consist of nonnegative nonincreasing profiles $h\leq1$,
constant with value $c$ on $[0,1)$, right continuous thereafter, satisfying
$\int_0^\infty h^3<\infty$. In particular $h(t)\to0$.
Set
\[
 \mathcal E^1_{\rm rad}=\{u_h+C:h\in\mathcal H,\ C\in\R\}.
\]
At $T=+\infty$ use the limit $g_h(+\infty)=\int_0^\infty h$, which may
be infinite. The only point of infinite depth is the rigid point $y$.
These functions are canonical decreasing limits of bounded generated
potentials; no recognition of an intrinsic unbounded psh class is asserted.

\begin{theorem}\label{thm:extension}
Equations \eqref{eq:bounded} and \eqref{eq:mixed} extend to all
$h,k\in\mathcal H$ by weak limits of actual bounded products. They define
positive analytic Radon measures of mass two, without an atom at $y$,
independently of the finite-step approximations described below. The
energy is
\[
 E(u_h+C)=-\frac16-\frac13\int_0^\infty h^3\,\mathrm dt+2C.
\]
Differences in measure integrals are taken off $y$ and assigned zero on
the null singleton $\{y\}$ if both potentials are infinite there.
For every $h,k\in\mathcal H$ the segment first variation is
\[
 \left.\frac{\mathrm d}{\mathrm d\epsilon}\right|_{0+}
 E(u_{(1-\epsilon)h+\epsilon k})
       =-\int_0^\infty h^2(k-h)\,\mathrm dt
       =\int_X(u_k-u_h)\,\MA(u_h).
\]
The products and energy are continuous along a decreasing sequence in
this generated class whose energies are bounded below. For mixed
products this requires such a sequence in each of its two slots.
\end{theorem}
\begin{proof}
Approximate $h$ from below by compact nonincreasing finite step functions
$h_m$, constant $c$ on $[0,1)$: use right-endpoint samples on nested
dyadic partitions of $[1,2^m)$ and set the tail to zero. At continuity
points $h_m\uparrow h$. Thus $g_{h_m}\uparrow g_h$ at every finite depth,
and at infinity by interchange of the two nonnegative suprema.
$u_{h_m}$ decreases to $u_h$, including at $y$. It is locally uniform
away from $y$ and its energy converges by dominated convergence of $h_m^3$.

The finite measures $-\mathrm d(h_m^2)$ converge on bounded depth
intervals, tested at continuity points of $h$, to the Stieltjes measure
$-\mathrm d(h^2)$. They are uniformly tight: their tail above a continuity
depth $R$ is at most $h(R)^2\to0$. Approximate a continuous test on a
bounded interval by its finite partition values, then discard this
uniformly small tail. This proves weak convergence against bounded
continuous depth tests. The same argument after composition with the
continuous actual section $t\mapsto\xi_t$ proves convergence against
analytic tests on $X$. It also proves independence for any compact finite
profiles converging almost everywhere with initial value tending to $c$
and the stated uniform tail bound. Their outside part converges by
\eqref{eq:bounded}. For two profiles use $h_mk_m$, whose tail tends
uniformly to zero; polarization proves \eqref{eq:mixed}. Its Stieltjes
mass is $cd$, while the outside mass is $2-cd$. This proves positivity
and full raw mass. There is no missing mass at depth infinity.

For integrability and variation, Stieltjes integration followed by
nonnegative layer cake gives
\[
 \int g_k(t)[-\mathrm d(h^2)](t)=\int_0^\infty k(t)h(t)^2\,\mathrm dt
                         \leq\|k\|_3\|h\|_3^2.
\]
Signed differences follow by subtraction. Their outside contribution
vanishes since $g_h(T)=0$ off $W$. Expanding the cubic energy along the
segment, with H\"older's inequality for all mixed terms, proves the
variation. The energy is the genuine limit of the bounded energy
formula, and agrees with the bounded operator whenever $g_h$ is bounded:
in that case the approximants converge uniformly on $X$, so the retained
bounded uniform continuity applies. Constants add raw mass times $C$.

We give the decreasing-continuity argument explicitly. The constant is
read from the value at any fixed maximum point of $u_0$ outside $F$,
where $u_h+C=u_0+C$. Along a decreasing sequence its constants $C_m$
decrease. If its energies are bounded below, then
$E(u_{h_m})\leq-1/6$ implies that $C_m$ is also bounded below.
Subtract their limit $C$; the functions
$g_m(t)-(C_m-C)$ increase pointwise to a finite concave limit at finite
depth, and their derivatives converge almost everywhere by the secant
inequalities for concave functions. Hence $h_m\to h$ almost everywhere.
Bounded energy and bounded constants bound $\int h_m^3$, so Fatou gives
$h\in\mathcal H$. The functions $g_m$ themselves converge locally
uniformly to $g_h$; the additive discrepancies tend to zero. At infinity
the increasing shifted functions give the same canonical limiting value.

For energy convergence with variable constants we need a little more
than Fatou. Put $\delta_m=C_m-C\downarrow0$. Decrease of the potentials
implies $g_m(t)\leq g_h(t)+\delta_m$ for every $t$. For a nonnegative
nonincreasing compact test $a(t)$, Stieltjes integration by parts gives
\[
 \int a h_m\leq\int a h+\delta_m a(0).
\]
Choose $a=h_m^2\mathbf1_{[0,R)}$. Then $a(0)\leq1$, and H\"older gives
$Z_m(R)\leq\|h\|_3 Z_m(R)^{2/3}+\delta_m$, where
$Z_m(R)=\int_0^R h_m^3$. Letting $R\to\infty$ and then taking limit
superiors shows $\limsup\int h_m^3\leq\int h^3$; Fatou gives the reverse
bound. Energies therefore converge. Monotonicity of the profiles and
the common cubic bound give $h_m(R)^2\leq(C'/R)^{2/3}$, uniformly.
This gives tightness and hence weak product convergence just as above;
in two slots use the product of the corresponding tail bounds.
\end{proof}

\section{Comparison and a separated compactness replacement}
\begin{theorem}\label{thm:comparison}
For any two generated radial potentials $u_h+C,u_k+D$,
\[
 \MA(u_k)\{u_h+C<u_k+D\}\leq
 \MA(u_h)\{u_h+C<u_k+D\}.
\]
For every $M<\infty$, profiles in $\mathcal H$ with $\int h^3\leq M$
form a compact set for local uniform convergence of $g_h$ on finite
depth intervals. This is a separated metrizable topology and is exactly
local uniform convergence of the canonical potentials away from $y$,
with the fixed normalization $\sup u_h=1/12$.
\end{theorem}
\begin{proof}
Put $f=g_h-g_k$ and $b=C-D$. The strict comparison set inside $W$
corresponds to the open depth set $\{f>b\}$. On each of its intervals
$(a,b')$ with finite contact endpoints, positivity just to the right of
$a$ gives $h(a+)\geq k(a+)$, while return to contact at $b'$ gives
$h(b'-)\leq k(b'-)$. These follow directly from the right and left
derivatives of the two concave primitives. Thus the difference of
Stieltjes masses on the interval is nonnegative:
\[
 [h(a+)^2-k(a+)^2]-[h(b'-)^2-k(b'-)^2]\geq0.
\]
At a right endpoint infinity both tails vanish. At a left endpoint zero,
if $b=0$ and the interval begins there, the same right derivative
inequality holds. If $b<0$, the initial derivative need not have that
sign, but the entire outside comparison set is included; its mass
difference is $d^2-c^2$, cancelling the left Stieltjes contribution
$c^2-d^2$. If $b>0$ no interval starts at zero and the outside set is
empty. Summing over the countably many intervals proves the inequality.
Endpoints where equality holds are excluded; the right/left tail
conventions in the displayed calculation account for their atoms.
The rigid point has zero measure in both slots.

For compactness, all $g_h$ vanish at zero, are nondecreasing and
concave, and are 1-Lipschitz. A diagonal subsequence of their values on
the nonnegative rationals converges; the Lipschitz bound gives uniform
convergence on every finite interval. The limit $g$ is concave with
derivative $h\in[0,1]$ almost everywhere, constant on $[0,1)$.
Concavity's secant inequalities imply almost everywhere derivative
convergence. Fatou gives $\int h^3\leq M$; take its right continuous
representative. Thus the limit belongs to $\mathcal H$. This proves
sequential compactness in the metric
$\sum_{m\geq1}2^{-m}\min(1,\sup_{[0,m]}|g_h-g_k|)$, and hence compactness.
Distinct primitives are distinguished on a finite interval and hence
by an actual Gauss point $\xi_t$; values at depths below one are fixed
by $g(1)$. The topology is separated.

The closed depth bidiscs $\{T\geq R\}$ decrease to $\{y\}$.
Every compact subset of $X\setminus\{y\}$ therefore has bounded depth,
so finite-interval convergence implies local uniform convergence there.
Conversely each Gauss depth segment $\{\xi_t:1\leq t\leq R\}$ is compact
and detects uniform primitive convergence, including the initial linear
piece. Finally $u_h\leq u_0$, with equality at its maximum off $F$, so
$\sup u_h=1/12$. This topology need not control the value at $y$; that
value is nevertheless uniquely fixed by the canonical profile.
\end{proof}

\section{The equation family and uniqueness}
\begin{theorem}\label{thm:equation}
Let $\nu$ be a positive Radon measure on $[1,\infty)$ of mass $c^2\leq1$
satisfying
\[
 \int_1^\infty\nu((t,\infty))^{3/2}\,\mathrm dt<\infty.
\]
Then the actual analytic measure
\[
 \mu_\nu=(1-c^2)\mu_0+c^2\lambda+\xi_*\nu
\]
has raw mass two and admits a solution in $\mathcal E^1_{\rm rad}$,
unique within that class modulo constants. It is given by
$h=c$ on $[0,1)$ and $h(t)=\sqrt{\nu((t,\infty))}$ for $t\geq1$.
This is precisely the image measure family of the constructed radial
operator; the tail condition is necessary for a radial finite-energy
solution in this family.
\end{theorem}
\begin{proof}
The stated profile is right continuous and nonincreasing, tends to zero,
and has cubic integral $c^3$ plus the displayed finite integral.
Its Stieltjes measure $-\mathrm d(h^2)$ equals $\nu$, including the jump
at 1. Theorem~\ref{thm:extension} therefore gives $\MA(u_h)=\mu_\nu$.
This is a genuine bounded-operator limit, not an operator defined by
assigning an inverse measure. Every finite-energy radial profile has this
tail condition, proving necessity and the image statement.
The map $t\mapsto\xi_t$ is continuous and injective; on each bounded
interval it is a homeomorphism onto its image, since $X$ is Hausdorff.
Thus equality of measures on $W$ recovers $\nu$ and all its tails.
Nonnegativity recovers $h$ uniquely almost everywhere and then its
primitive, its canonical value at $y$, and the normalized potential.
Only the free additive constant remains.
\end{proof}

\section{Scope, source dependence, and the remaining original problem}
These are complete conditional proofs for a generated branch class.
The measure family fixes its outside part and restricts its inside part
to one actual translated Gauss branch. Uniqueness is only within the
radial class; compactness, variation, and comparison are only as stated
there. Finite dual energy for a larger class is not established for every
$\mu_\nu$. No all-direction envelope or orthogonality theorem is supplied.
In particular this paper does not close the intrinsic all-direction
finite-energy and EVERY admissible arbitrary-support requirements of
P16.7. All original counts remain unchanged.

The exact bounded foundations and nonalgebraizable quotient geometry
are the unreviewed inputs of continuation 383, enumerated in continuation
384. The entire 384 proof was read; its source's CLD-v2 Sections 5.8.1--6
and immediate integration proofs were freshly read as documented there.
No new full-primary, source-pixel, recursive foundation, or entire-book
reading is claimed in this extension. All new Stieltjes, tail-tightness,
majorization, comparison, compactness and equation proofs are given here.
All full primary bodies remain local; no supplied or native program was
executed or imported. Independent mathematical review and compilation
verification are pending. The existing P16.4 editor remains unchanged.
\end{document}
